% Pursuing Stacks (cotes 134-2 to 134-8): archive cote 134-5, site page 100 (PDF page 101), typescript page 287.
% typescript. Draft compared: ps4.tex:5123-5202.
% AI-transcribed by Claude Opus 5.5 (claude-opus-5-5, adaptive thinking, effort high), Message Batches API, batch ps2 (msgbatch_01NVwr1Yuo8qJf5rykeGBXPK), 2026-11-12; reviewed against both supplied scan halves; see CHECK_DECISIONS.tsv.
\scanpage{134-5/100}
\mspage{287}

Stated this way, the condition just obtained on the pair of adjoint
functors $(f_!,f^*)$\hand{,} between the two \xout{contractibility structures (} categories
M, M' endowed with contractibility structures, namely (2) and
\[ (4) \qquad f^*(M_a) \hand{\subset} M'_a \quad , \]
are not quite symmetric, as the condition on $f_!$ is expressed in terms
of contractible objects, whereas the condition on $f^*$ is expressed in
terms of aspheric \deleted{objects} \inserted{ones}. % One continuous oblique pen stroke runs from assuming/that toward the gap after (4) on the next line.
However, assuming\hand{/}that $\underline{W}$ satisfies Loc 4)
and using theorem 1 p. 252, we see that (4) is equivalent to the
condition
\[ (5) \qquad f^*(M_c) \hand{\subset} M'_c \quad , \]
which does not depend any longer on the choice of $\underline{W}$ !

Finally, we are led to a notion of pure homotopy theory, in
terms of contractibility structures alone, without the intr\xout{o}usion of
the choice of a basic localizer $\underline{W}$ and corresponding asphericity notions\hand{.}
We may call \xout{it} the pair a ``\hand{\uline{\textcolor{black}{bimorphism}}}'' of contractibility structures,
and introduce it via the following \hand{summing-up statement:}

\inserted{``Scholie''} \deleted{\hand{\uline{\textcolor{black}{Proposition}}}} : Let $(M,M_c)$ and $(M',M'_c)$ be two contractibility structure,
each admitting a small generating subcategory for the contractibility
structures, say C and C' respectively. Let $(f_!,f^*)$ be a pair of ad-%
joint functors
\[ f_! : M' \hand{\longrightarrow} M \quad , \qquad f^* : M \hand{\longrightarrow} M' \quad . \]
Let's consider the following inclusion conditions
\[ (!) \ \xout{(!\illegible{1}} \qquad f_!(M'_c) \hand{\subset} M_c \quad , \]
\margin{(!')}
\[ (!\hand{'}) \qquad f_!(C') \hand{\subset} M_o \quad , \]
\[ (*) \qquad f^*(M_c) \hand{\subset} M'_c \quad , \]
\[ (*') \qquad f^*(C) \hand{\subset} M'_o \quad , \]
where $M_o$ and $M'_o$ are (as in section 81) the sets of ``0-connected''
objects in M and M' respectively, for the given contractibility struc-%
tures.
% A pen step/caret marks the new paragraph boundary before a).
\par\hand{a)} \deleted{(}Clearly, in view of
\[ C \hand{\subset} M_c \hand{\subset} M_o \quad , \quad C' \hand{\subset} M'_c \hand{\subset} M'_o \quad , \]
condition (!) implies (!') and condition (*) implies (*')\deleted{)}. Moreover,
it is even true that conditions (*) and (*') are equivalent, and the
same holds for (!) and (!') provided $f_!$ commutes with finite pro\xout{\unsure{ducts}}\hand{-}%
ducts. We'll say that $(f_!\ ,\ f^*\ )$ is a \hand{\uline{\textcolor{black}{bimorphism}}} for the given contrac-%
tibility structures, if \inserted{the} inclusions (!) and (*) do hold.

\margin{100}
